\documentclass{replab}
\usepackage{amsfonts}
\usepackage{amssymb}
\usepackage{amsthm}
\usepackage{array}
\usepackage[english, spanish, es-noitemize, es-noindentfirst]{babel}
\usepackage{bm}
\usepackage{booktabs}
\usepackage{cancel}
\usepackage[nolimits]{cmupint}
\usepackage{color}
\usepackage{empheq}
\usepackage{fix-cm}
\usepackage{lipsum}
\usepackage{listing}
\usepackage[bb=ams]{mathalpha}
\usepackage{mathptmx}
\usepackage{mathrsfs}
\usepackage{mathtools}
\usepackage{multicol}
\usepackage{multirow}
\usepackage{parskip}
\usepackage{pgfplotstable}
\usepackage{psfrag}
\usepackage{siunitx}
\usepackage{tabularx}
\usepackage{tikz}
\usepackage{verbatim}
\usepackage{wrapfig}
\usepackage{xcolor}
\usepackage{xfrac}

\hypersetup{
    pdftitle={Diferencias Finitas - Ecuación de Onda 1D}
    colorlinks=true,
    linkcolor=prussian,
    citecolor=prussian,
    filecolor=prussian,
    urlcolor=prussian,
    }
    
% --- Información del documento ---
\title{Background}
\subtitle={Diferencias Finitas - Ecuación de Onda 1D}
\author{Santiago Talero Parra}
\date{\today}
\subject={Modelamiento Físico Computacional}
\email={\href{mailto:stalerop@udistrital.edu.co}{\texttt{stalerop@udistrital.edu.co}}}

% --- Inicio del documento ---
\begin{document}
\pagestyle{fancy}
\unspacedoperators
\begin{center}
    \maketitle{\hrulefill}
\end{center}
\selectlanguage{spanish}

% --- Cuerpo del reporte ---

\section*{Ecuación de ondas}
Para ilustrar como se hallar una solución aproximada de una EDP hiperbólica, vamos a considerar la denominada ecuación de ondas:
\begin{equation*}
	\frac{\partial^2}{\partial t^2}u(x,t) - c^2 \frac{\partial^2}{\partial x^2}u(x,t) = 0, \quad x \in [0,L], \quad t > 0,
\end{equation*}

que modela la vibración de una cuerda elástica de longitud $L$ a lo largo del tiempo $t$ y $\alpha$ es un parámetro que depende de la cuerda considerada. EL valor de $c$ depende de las propiedades físicas del medio en el que se propaga la onda. En este caso, $u(x,t)$ representa el desplazamiento vertical de la cuerda en la posición $x$ y en el tiempo $t$. La ecuación de ondas es una EDP de segundo orden en el tiempo y en el espacio. Para resolver este problema, es necesario especificar las condiciones iniciales y de frontera. Las condiciones iniciales son:
\begin{align*}
	u(x,0) &= f(x), \quad x \in [0,L], \\
	\frac{\partial}{\partial t}u(x,0) &= g(x), \quad x \in [0,L],
\end{align*}
donde $f(x)$ y $g(x)$ son funciones dadas que representan la posición inicial y la velocidad inicial de la cuerda, respectivamente. Las condiciones de frontera son:
\begin{align*}
	u(0,t) &= 0, \quad t \geq 0, \\
	u(L,t) &= 0, \quad t \geq 0,
\end{align*}
lo que significa que los extremos de la cuerda están fijos en el tiempo.

\section*{Malla}

Seleccionamos el paso espacial o para la variable $x$ como un entero $m > 0$ y definimos $h = \frac{L}{m}$. Luego, definimos los nodos espaciales como $x_i = ih$ para $i = 0, 1, \ldots, m$. De manera similar, seleccionamos el paso temporal o para la variable $t$ como un entero $k > 0$ y definimos $t_j = jk$ para $j = 0, 1, \ldots$. El objetivo es hallar una aproximación $\hat{u}(x_i,t_j), \quad i=0,1, \ldots, m \quad j=0,1, \ldots$ en los puntos del mallado. 

\section*{Esquema de diferencias finitas}
Para aproximar las derivadas parciales en la ecuación de ondas, utilizamos diferencias finitas centradas. La segunda derivada temporal se aproxima como:
\begin{equation*}
	\frac{\partial^2}{\partial t^2}u(x_i,t_j) \approx \frac{u(x_i,t_{j+1}) - 2u(x_i,t_j) + u(x_i,t_{j-1})}{k^2}
\end{equation*}
y la segunda derivada espacial se aproxima como:
\begin{equation*}
	\frac{\partial^2}{\partial x^2}u(x_i,t_j) \approx \frac{u(x_{i+1},t_j) - 2u(x_i,t_j) + u(x_{i-1},t_j)}{h^2}
\end{equation*}
Sustituyendo estas aproximaciones en la ecuación de ondas, obtenemos:
\begin{equation*}
	\frac{u(x_i,t_{j+1}) - 2u(x_i,t_j) + u(x_i,t_{j-1})}{k^2} - c^2 \frac{u(x_{i+1},t_j) - 2u(x_i,t_j) + u(x_{i-1},t_j)}{h^2} = 0
\end{equation*}
Reorganizando esta ecuación, obtenemos el esquema de diferencias finitas explícito:
\begin{equation}
	u(x_i,t_{j+1}) = 2u(x_i,t_j) - u(x_i,t_{j-1}) + \lambda^2 \left( u(x_{i+1},t_j) - 2u(x_i,t_j) + u(x_{i-1},t_j) \right)
	\label{eqn1}
\end{equation}
donde $\lambda = \frac{\alpha k}{h}$ es el número de Courant.\\

El sistema de ecuaciones que verifica el vector de aproximaciones $\hat{u}^{(j+1)}$ está dado por la ecuación \ref{eqn1} para $i=1, 2, \ldots, m-1$ y las condiciones de frontera $\hat{u}^{(j+1)}_0 = 0$ y $\hat{u}^{(j+1)}_m = 0$.\\

Matricializando el sistema, obtenemos la siguiente relación de recurrencia:
\begin{equation*}
	\hat{u}^{(j+1)} = (2 - 2\lambda^2) \hat{u}^{(j)} + \lambda^2 \mathbf{A} \hat{u}^{(j)} - \hat{u}^{(j-1)}
\end{equation*}
donde $\mathbf{A}$ es la matriz tridiagonal dada por:
\begin{equation*}
\mathbf{A} = 
\begin{pmatrix}
0 & 1 & 0 & \cdots & 0 \\
1 & 0 & 1 & \ddots & \vdots \\
0 & 1 & 0 & \ddots & 0 \\
\vdots & \ddots & \ddots & \ddots & 1 \\
0 & \cdots & 0 & 1 & 0
\end{pmatrix}_{(m-1)\times(m-1)}
\end{equation*}
Esta matriz representa el operador de diferencias finitas para la segunda derivada espacial, considerando las condiciones de frontera homogéneas.

Despejando $u_{i, j+1}$ (el valor en el nuevo nivel temporal):
$$ u_{i, j+1} = 2u_{i, j} - u_{i, j-1} + \lambda^2 (u_{i+1, j} - 2u_{i, j} + u_{i-1, j}) $$
$$ u_{i, j+1} = 2(1 - \lambda^2)u_{i, j} + \lambda^2(u_{i+1, j} + u_{i-1, j}) - u_{i, j-1} \quad (*) $$

\subsection*{Representación 1: La Matriz Completa ($\mathbf{M}$)}

Esta representación escribe la relación de recurrencia directamente como una operación lineal sobre el vector de soluciones $\mathbf{\hat{u}}^{(j)}$ en el tiempo $t_j$:
$$ \mathbf{\hat{u}}^{(j+1)} = \mathbf{M} \mathbf{\hat{u}}^{(j)} - \mathbf{\hat{u}}^{(j-1)} $$
donde $\mathbf{\hat{u}}^{(j)} = [u_{1, j}, u_{2, j}, \dots, u_{m-1, j}]^T$ es el vector de desplazamientos en el tiempo $t_j$.
La matriz $\mathbf{M}$ es una matriz tridiagonal de tamaño $(m-1) \times (m-1)$, definida por los coeficientes de $u_{\cdot, j}$ en la ecuación $(*)$:
$$ \mathbf{M} = 
\begin{pmatrix}
2(1 - \lambda^2) & \lambda^2 & 0 & \cdots & 0 \\
\lambda^2 & 2(1 - \lambda^2) & \lambda^2 & \ddots & \vdots \\
0 & \lambda^2 & 2(1 - \lambda^2) & \ddots & 0 \\
\vdots & \ddots & \ddots & \ddots & \lambda^2 \\
0 & \cdots & 0 & \lambda^2 & 2(1 - \lambda^2)
\end{pmatrix} $$

\subsection*{Representación 2: El Operador Espacial Separado ($\mathbf{A}$)}

La segunda representación, que posee la matriz tridiagonal $\mathbf{A}$, factoriza la matriz en términos de la identidad ($\mathbf{I}$) y la matriz tridiagonal ($\mathbf{A}$) que representa la suma de los vecinos $(u_{i+1, j} + u_{i-1, j})$:
$$ \mathbf{\hat{u}}^{(j+1)} = (2 - 2\lambda^2) \mathbf{\hat{u}}^{(j)} + \lambda^2 \mathbf{A} \mathbf{\hat{u}}^{(j)} - \mathbf{\hat{u}}^{(j-1)} \quad (**) $$
donde $\mathbf{A}$ es:
$$ \mathbf{A} = 
\begin{pmatrix}
0 & 1 & 0 & \cdots & 0 \\
1 & 0 & 1 & \ddots & \vdots \\
0 & 1 & 0 & \ddots & 0 \\
\vdots & \ddots & \ddots & \ddots & 1 \\
0 & \cdots & 0 & 1 & 0
\end{pmatrix} $$

\subsection*{Demostración de Equivalencia}

Para demostrar la equivalencia, debemos mostrar que el factor que multiplica a $\mathbf{\hat{u}}^{(j)}$ en $(**)$ es igual a $\mathbf{M}$:
$$ \mathbf{M} \stackrel{?}{=} (2 - 2\lambda^2) \mathbf{I} + \lambda^2 \mathbf{A} $$
Donde $\mathbf{I}$ es la matriz identidad.

$$ (2 - 2\lambda^2) \mathbf{I} + \lambda^2 \mathbf{A} = (2(1-\lambda^2))
\begin{pmatrix}
1 & 0 & 0 & \cdots \\
0 & 1 & 0 & \cdots \\
0 & 0 & 1 & \cdots \\
\vdots & \vdots & \vdots & \ddots
\end{pmatrix} 
+ \lambda^2
\begin{pmatrix}
0 & 1 & 0 & \cdots \\
1 & 0 & 1 & \cdots \\
0 & 1 & 0 & \cdots \\
\vdots & \vdots & \vdots & \ddots
\end{pmatrix} $$

Aplicando la suma matricial:

\begin{itemize}
    \item \textbf{Diagonal Principal ($i=k$):} El elemento $M_{i,i}$ es el resultado de la suma de los términos diagonales:
    $$ M_{i,i} = (2(1-\lambda^2)) \cdot 1 + \lambda^2 \cdot 0 = 2(1 - \lambda^2) $$
    \item \textbf{Sub-diagonales ($|i-k|=1$):} El elemento $M_{i,i\pm 1}$ es el resultado de la suma de los términos fuera de la diagonal:
    $$ M_{i,i\pm 1} = (2(1-\lambda^2)) \cdot 0 + \lambda^2 \cdot 1 = \lambda^2 $$
\end{itemize}

Sustituyendo estos valores, obtenemos:
$$ (2 - 2\lambda^2) \mathbf{I} + \lambda^2 \mathbf{A} = 
\begin{pmatrix}
2(1 - \lambda^2) & \lambda^2 & 0 & \cdots & 0 \\
\lambda^2 & 2(1 - \lambda^2) & \lambda^2 & \ddots & \vdots \\
0 & \lambda^2 & 2(1 - \lambda^2) & \ddots & 0 \\
\vdots & \ddots & \ddots & \ddots & \lambda^2 \\
0 & \cdots & 0 & \lambda^2 & 2(1 - \lambda^2)
\end{pmatrix} = \mathbf{M} $$

El sistema anterior tiene un problema: cuando intentamos calcular el vector de aproximaciones $\hat{u}^{(1)}$ para $t=t_1$ que resulta de aplicar la recurrencia anterior para $j=0$, necesitamos conocer $\hat{u}^{(-1)}$ ¡Cuyos valores no conocemos!\\

Observemos que aún no hemos aplicado la restricción de la condición inicial de velocidad, es decir, la segunda condición inicial:
$$ \frac{\partial}{\partial t}u(x,0) = g(x), \quad x \in [0,L] $$
que en términos de las aproximaciones en los nodos espaciales es:
$$ \frac{\partial}{\partial t}u(x_i,0) = g(x_i), \quad i=0,1, \ldots, m $$
Podemos utilizar esta condición para obtener una expresión para $\hat{u}^{(1)}$ en función de $\hat{u}^{(0)}$ y $g(x_i)$. Dicha restricción nos permitirá conocer el vector de aproximaciones $\hat{u}^{(1)}$ de la forma siguiente:
Aproximamos $\frac{\partial}{\partial t}u(x_i,0)$ usando la diferencia dividida siguiente:
\begin{align*}
\frac{\partial g}{\partial t}(x_i,0)= & \frac{u(x_i,t_1)-u(x_i,0)}{k}-\frac{k}{2}\frac{\partial^2 u}{\partial t^2}(x_i,\tilde{\mu}_i)\\ = & \frac{u(x_i,t_1)-f(x_i)}{k}-\frac{k}{2}\frac{\partial^2 u}{\partial t^2}(x_i,\tilde{\mu}_i)=g(x_i),
\end{align*}

donde $\tilde{\mu}_i \in (0,t_1)$. Despejando $u(x_i,t_1)$, obtenemos:
$$ u(x_i,t_1) = f(x_i) + k g(x_i) + \frac{k^2}{2} \frac{\partial^2 u}{\partial t^2}(x_i,\tilde{\mu}_i) $$
Aproximando $\frac{\partial^2 u}{\partial t^2}(x_i,\tilde{\mu}_i)$ usando la ecuación de ondas, obtenemos:
$$ \frac{\partial^2 u}{\partial t^2}(x_i,\tilde{\mu}_i) = c^2 \frac{\partial^2 u}{\partial x^2}(x_i,\tilde{\mu}_i) $$
Aproximando $\frac{\partial^2 u}{\partial x^2}(x_i,\tilde{\mu}_i)$ usando diferencias finitas centradas, obtenemos:
$$ \frac{\partial^2 u}{\partial x^2}(x_i,\tilde{\mu}_i) \approx \frac{u(x_{i+1},\tilde{\mu}_i) - 2u(x_i,\tilde{\mu}_i) + u(x_{i-1},\tilde{\mu}_i)}{h^2} $$
Sustituyendo esta aproximación en la expresión para $u(x_i,t_1)$, obtenemos:
$$ u(x_i,t_1) \approx f(x_i) + k g(x_i) + \frac{c^2 k^2}{2h^2} \left( u(x_{i+1},\tilde{\mu}_i) - 2u(x_i,\tilde{\mu}_i) + u(x_{i-1},\tilde{\mu}_i) \right) $$

El problema es que hemos usado una aproximación con error de orden lineal temporal $O(k)$ y cuadrático espacial $O(h^2)$ para obtener una expresión para $u(x_i,t_1)$. Sin embargo, podemos mejorar la precisión de esta expresión utilizando el esquema de diferencias finitas explícito \ref{eqn1} para aproximar $u(x_i,t_1)$. Esto hará que nuestras aproximaciones futuras se vean afectadas por dicho error.\\

Para mejorar la aproximación de $u(x_i,t_1)$, desarrollamos por Taylor $u(x_i,t_1)$ alrededor de $t=0$ hasta el segundo orden para $i=1,2,\ldots,m-1$:

\begin{align*}
u(x_i,t_1)=& u(x_i,0)+k\frac{\partial u}{\partial t}(x_i,0)+\frac{k^2}{2}\frac{\partial^2 u}{\partial t^2}(x_i,0)+\frac{k^3}{6}\frac{\partial^3 u}{\partial t^3}(x_i,\hat{\mu}_i)\\ 
= & f(x_i)+k g(x_i)+\frac{k^2}{2}\frac{\partial^2 u}{\partial t^2}(x_i,0)+\frac{k^3}{6}\frac{\partial^3 u}{\partial t^3}(x_i,\hat{\mu}_i),
\end{align*}

donde $\hat{\mu}_i \in (0,t_1)$. Usando la ecuación de ondas, podemos escribir $\frac{\partial^2 u}{\partial t^2}(x_i,0)$ al considerar que $u(x_i,0)=f(x_i)$, si suponemos que $f$ es dos veces diferenciable, obtenemos:
$$ \frac{\partial^2 u}{\partial t^2}(x_i,0) = c^2 \frac{\partial^2 f}{\partial x^2}(x_i) $$
Aproximando $\frac{\partial^2 f}{\partial x^2}(x_i)$ usando diferencias finitas centradas, obtenemos:
$$ \frac{\partial^2 f}{\partial x^2}(x_i) \approx \frac{f(x_{i+1}) - 2f(x_i) + f(x_{i-1})}{h^2} $$
Sustituyendo esta aproximación en la expresión para $u(x_i,t_1)$, obtenemos:
$$ u(x_i,t_1) \approx f(x_i) + k g(x_i) + \frac{c^2 k^2}{2h^2} \left( f(x_{i+1}) - 2f(x_i) + f(x_{i-1}) \right) + \frac{k^3}{6} \frac{\partial^3 u}{\partial t^3}(x_i,\hat{\mu}_i) $$
Esta expresión tiene un error de orden $O(k^3 + k^2 h^2)$, que es más preciso que la aproximación anterior.\\

En el caso de que no exista la derivada segunda de la función $f$, podemos aproximar $f''(x_i)$ usando diferencias finitas centradas directamente sobre los valores de $f$ en los nodos espaciales:
\[f''(x_i)=\frac{f(x_{i+1})-2f(x_i)+f(x_{i-1})}{h^2}-\frac{h^2}{12}f^{(4)}(\tilde{\xi}),\]
donde $\tilde{\xi} \in (x_{i-1},x_{i+1})$. Así, la expresión para $u(x_i,t_1)$ queda como:
\begin{align*}
\hat{u}_{i1}= & \hat{u}(x_i,t_1)=f(x_i)+kg(x_i)+\frac{\alpha^2 k^2}{2h^2}(f(x_{i+1})-2f(x_i)+f(x_{i-1})),\\
= & f(x_i)+kg(x_i)+\frac{\lambda^2}{2}(f(x_{i+1})-2f(x_i)+f(x_{i-1}))\\
= & \left(1-\lambda^2\right)f(x_i)+\frac{\lambda^2}{2}(f(x_{i+1})+f(x_{i-1}))+kg(x_i),
\end{align*}
para $i=1,2,\ldots,m-1$ con error de truncamiento $O(k^3+h^2)$.\\








\printbibliography[heading=bibintoc]

\end{document}